
When AI assistants match user communication style, do users engage more? Intuition suggests that accommodation---adapting linguistic style to match an interlocutor---should foster rapport and encourage continued interaction. Communication Accommodation Theory (CAT) predicts that convergence signals attentiveness and understanding \citep{giles1973accommodation}. Yet this prediction, grounded in human-human interaction research, may not transfer directly to human-AI dialogue.

This paper presents an empirical analysis of formality accommodation in 111,039 human-AI conversations, testing whether accommodation predicts engagement as measured by conversation continuation. The central finding challenges linear assumptions: moderate accommodation predicts highest engagement, while both under- and over-accommodation are associated with reduced continuation rates.

\subsubsection{1.1 The Problem of Optimal Accommodation}

Prior work establishes that linguistic accommodation exists in human-AI interaction. Chen et al. (2026) demonstrated bidirectional accommodation in GPT-4o conversations, finding that model adaptation is front-loaded while user convergence is gradual. However, this work measured whether accommodation occurs, not how much accommodation is associated with optimal engagement. The relationship between accommodation magnitude and interaction outcomes has not been tested at scale.

This gap has practical implications for AI design. If maximal accommodation improves engagement, chatbots should mirror user style as closely as possible. If accommodation has diminishing returns---or negative effects at extremes---designers need to calibrate adaptation carefully.

\subsubsection{1.2 Contributions}

This study addresses this gap through systematic analysis of formality accommodation in the Anthropic hh-rlhf dataset. The analysis proceeds through four validated sub-hypotheses:

\begin{enumerate}
\item \textbf{Accommodation is measurable} (H-E1): Bidirectional Convergence Score variance is non-trivial (SD=0.569, n=26,395), confirming that accommodation signals exist in the data.
\end{enumerate}

\begin{enumerate}
\item \textbf{AI adapts to human formality} (H-M2): AI response formality correlates with human input formality (r=0.152, p<0.001, n=111,039).
\end{enumerate}

\begin{enumerate}
\item \textbf{Accommodation is asymmetric} (H-M1, H-M2): AI-to-human accommodation is approximately 11 times stronger than human-to-AI adaptation.
\end{enumerate}

\begin{enumerate}
\item \textbf{Optimal accommodation is non-linear} (H-M3): Conversations in the middle tercile of formality delta show 71.4\% continuation rate, compared to 65.9\% for low-delta and 60.9\% for high-delta conversations.
\end{enumerate}

The fourth finding is the primary contribution. While the original prediction assumed monotonic benefits from convergence, the data reveal an inverted-U relationship.

